\begin{lemma}[Wild odd upper residual normalization]
\label{mod:cases-wildodd-upperresidualnormalization}
Let $t=s+1$ be the last lower break.  For the coherent High source put
$A=N(\alpha _1)/\gamma_F$ and $u=\beta _1/\alpha _1$; for the actual Low
stationary norm pair put $A=\alpha/\delta$ and
$u=\epsilon _1\beta _1/\alpha _1$.  In both cases $A$ is retained as a unit
and has order
\[
 -\bigl((t+1)+n(\psi_F)\bigr).
\]
The literal High representative and every explicitly supplied Low upstairs
witness, divided by their respective unchanged denominators, are exactly
$A\bigl(N(u)-u\bigr)$.  Thus the source normalization uses the same
coherent High product/norm rows and the same Low stationary pair, parameter
table, and witness as PhaseReduction.

For the last-break displacement $x_z=[z]\pi_K^t$, its norm polynomial lies
on the exact layer $t$.  Specializing the real High and Low raw-factor
theorems to $x_z$ gives
\[
 \psi_F\bigl(A P(x_z)\bigr)=1.
\]
Here the indexed norm character is removed only because the displayed unit
is exhibited in the norm range.  If $\lambda\ne0$ is the ramification
coordinate, the generic positive-break critical-norm theorem from
Theorem~\ref{mod:ramification-normbelowbreak}, specialized at $t=s+1$,
gives exactly
\[
 \overline P(z)=z^p-\lambda^{p-1}z,
\]
and the raw identity gives the exact annihilator
\[
 \psi_0\!\left(\eta_0
   (z^p-\lambda^{p-1}z)\right)=1.
\]
Consequently the inverse Frobenius root has the manuscript's orientation
\[
 \operatorname{Frob}^{-1}(\eta_0)=\eta_0\lambda^{p-1}.
\]

For $y$ of order $r$ with $r+2d_K=t$, and with its normalized initial
residue equal to $-\iota$, the common-coordinate trace specialization is
proved exactly as
\[
 \psi_F\!\left(A\operatorname{Tr}_{K/F}
   (y\pi_K^{2d_K}[z])\right)
 =\psi_0\bigl((\beta^{-1}\iota)z\bigr),
 \qquad
 \beta^{-1}=\eta_0\lambda^{p-1}.
\]
This is the precise \texttt{hnormalizedTrace} statement: its signed depth,
initial-residue hypothesis, trace direction, and lower residual coordinate
are explicit.

At the Low noncancelling boundary, where the actual conductor is $t+1$, the
normalized ratio has order zero.  Both it and its norm are integral.  For the
literal integral lift $n_O$ of $N(u)$ and
$d_0=\zeta^p$ in the common lower residual coordinate, one has
\[
  \operatorname{residueMap}(n_O)=d_0,
  \qquad d_0=\zeta^p,
\]
with $\zeta\ne0$ and $d_0\ne0$.  All denominator nonvanishing follows from
the unit-valued coefficients.  These statements are packaged in the narrow
normalization API consumed by UpperResidualCoefficients.

This node contains no upper polar rows, upper affine rows,
HigherSymmetricVanishing argument, final coefficient table, or parity
cancellation.
\LeanFile{LanglandsFirstMainLemma/Cases/WildOdd/UpperResidualNormalization.lean}
\lean{LanglandsFirstMainLemma.wildOdd_upperResidualNormalization}
\lean{LanglandsFirstMainLemma.wildOddHighStationaryRatio_eq_normalized}
\lean{LanglandsFirstMainLemma.wildOddLowStationaryRatio_eq_normalized}
\lean{LanglandsFirstMainLemma.wildOddCompactNormalizedTraceCharacter}
\lean{LanglandsFirstMainLemma.wildOddHighNormalization_raw}
\lean{LanglandsFirstMainLemma.wildOddLowNormalization_raw}
\lean{LanglandsFirstMainLemma.wildOddHighNormalization_annihilator}
\lean{LanglandsFirstMainLemma.wildOddLowNormalization_annihilator}
\lean{LanglandsFirstMainLemma.wildOddLowBoundaryNormNormalization_spec}
\leanok
\SourceText{Lemma lem:odd-break-line-normalization and the normalization and initial-coordinate portions of Lemma lem:residual-parameter-transport.}
\uses{mod:parameters-phasereduction,mod:ramification-normbelowbreak}
\FormalizationContract
The certificate exposes only the two actual source normalizations, the exact
last-break trace and norm-polynomial coordinate, the raw and annihilator
identities, inverse-Frobenius direction, and the integral boundary norm.
Coefficient rows, higher-symmetric vanishing, and finite-field cancellation
remain in their separate nodes.
\end{lemma}
